%; whizzy paragraph -pdf xpdf -latex ./whizzypdfptex.sh
%; whizzy-paragraph "^\\\\begin{frame}"
% latex beamer presentation.
% platex, latex-beamer でコンパイルすることを想定。 

%     Tokyo Debian Meeting resources
%     Copyright (C) 2011 Junichi Uekawa
%     Copyright (C) 2011 Nobuhiro Iwamatsu

%     This program is free software; you can redistribute it and/or modify
%     it under the terms of the GNU General Public License as published by
%     the Free Software Foundation; either version 2 of the License, or
%     (at your option) any later version.

%     This program is distributed in the hope that it will be useful,
%     but WITHOUT ANY WARRANTY; without even the implied warreanty of
%     MERCHANTABILITY or FITNESS FOR A PARTICULAR PURPOSE.  See the
%     GNU General Public License for more details.

%     You should have received a copy of the GNU General Public License
%     along with this program; if not, write to the Free Software
%     Foundation, Inc., 51 Franklin St, Fifth Floor, Boston, MA  02110-1301 USA

\documentclass[cjk,dvipdfmx,12pt]{beamer}
\usetheme{Tokyo}
\usepackage{monthlypresentation}

%  preview (shell-command (concat "evince " (replace-regexp-in-string "tex$" "pdf"(buffer-file-name)) "&")) 
%  presentation (shell-command (concat "xpdf -fullscreen " (replace-regexp-in-string "tex$" "pdf"(buffer-file-name)) "&"))
%  presentation (shell-command (concat "evince " (replace-regexp-in-string "tex$" "pdf"(buffer-file-name)) "&"))

%http://www.naney.org/diki/dk/hyperref.html
%日本語EUC系環境の時
\AtBeginDvi{\special{pdf:tounicode EUC-UCS2}}
%シフトJIS系環境の時
%\AtBeginDvi{\special{pdf:tounicode 90ms-RKSJ-UCS2}}

\title{東京エリアDebian勉強会}
\subtitle{第76回 2011年6月度}
\author{岩松 信洋 iwamatsu@debian.org\\IRC nick: iwamatsu}
\date{2011年6月18日}
\logo{\includegraphics[width=8cm]{image200607/openlogo-light.eps}}

\begin{document}

\frame{\titlepage{}}

\section{}
\begin{frame}
 \frametitle{Agenda}
\begin{minipage}[t]{0.45\hsize}
  \begin{itemize}
  \item 注意事項
	\begin{itemize}
	 \item 飲酒禁止
	 \item 宗教禁止
	 \item 営利活動禁止
	\end{itemize}
   \item 最近あったDebian関連のイベント報告
	\begin{itemize}
	 \item 第76回 東京エリア Debian勉強会
	\end{itemize}
 \end{itemize}
\end{minipage} 
\begin{minipage}[t]{0.45\hsize}
 \begin{itemize}
  \item Debian JP 定例会議処理系にXSLTを使ってみた
  \item Debian で sphinx と doxygen を使う
  \item OSC仙台 報告
 \end{itemize}
\end{minipage}
\end{frame}

\begin{frame}
 \frametitle{前回}
\begin{minipage}[t]{0.45\hsize}
  \begin{itemize}
  \item 注意事項
	\begin{itemize}
	 \item 飲食禁止
	 \item 宗教禁止
	 \item 営利活動禁止
	\end{itemize}
  \end{itemize}
\end{minipage}
\begin{minipage}[t]{0.45\hsize}
\begin{itemize}
  \item 最近あったDebian関連のイベント報告
 \begin{itemize}
  \item Apache2 のモジュールをつくってみた
  \item Debian on NiftyCloud
  \item Debian/m68k 開発
  \item 月刊PPC64ポーティング
 \end{itemize}
 \end{itemize}
\end{minipage}
\end{frame}


\emtext{イベント報告}

\emtext{DWN quiz}

\section{DWN quiz}
\begin{frame}{Debian 常識クイズ}

Debian の常識、もちろん知ってますよね?
知らないなんて恥ずかしくて、知らないとは言えないあんなことやこんなこと、
みんなで確認してみましょう。

今回の出題範囲は\url{debian-devel-announce@lists.debian.org} に投稿された
内容とDebian Project Newsからです。

\end{frame}

\subsection{問題}
%; whizzy-master ../debianmeetingresume201101.tex
% $B0J>e$N@_Dj$r$7$F$$$k$?$a!"$3$N%U%!%$%k$G(B M-x whizzytex $B$9$k$H!"(Bwhizzytex$B$,MxMQ$G$-$^$9!#(B
%
% $B$A$J$_$K!"%/%$%:$OJL%V%i%s%A$G:n@.$7!"$N$A$K%^!<%8$7$^$9!#5U$K%^!<%8$7(B
% $B$J$$$h$&$K$7$^$7$g$&!#(B
% (shell-command "git checkout quiz-prepare")

\santaku
{alioth.debian.org$B$,(B2$BBf$KJ,$+$l$^$7$?!#$=$N%5!<%PL>$O!)(B}
{vasks.debian.org $B$H(B wagner.debian.org}
{volks.debian.org $B$H(B don.debian.org}
{dennys.debian.org $B$H(B gusto.debian.org}
{A}
{$B$[$+$O%U%!%_%l%9$NL>A0(B}

\santaku
{$B8=:_9T$o$l$F$$$k(BPerl transition $B$N(BPerl$B%P!<%8%g%s$O!)(B}
{5.12}
{5.13}
{5.14}
{A}
{5.14$B$O$^$@(Bexperimental$B$G$9!#(B}

\santaku
{$B%W%i%$%^%j%_%i!<%5!<%P$,?7$7$/DI2C$5$l$?9q$O!)(B}
{$B%A%e%K%8%"(B}
{$BCf9q(B}
{$B%^%@%,%9%+%k(B}
{B}
{$B%A%e%K%8%"$H%^%@%+%9%+%k$O%_%i!<!#%W%i%$%^%j$G$O$J$$!#(B}


\emtext{prework}

{\footnotesize
 
\begin{prework}{ やまだ }

最近新しいカーネルを自分でビルドすることが再び
増えてきたので、その周りでごそごそやってます。
\begin{enumerate}
\item gitでタグを打って
\item そのタグをバージョンにして
\item ビルドして
\item テストVMに導入して
\item ついでに配布サーバに設置
\end{enumerate}

の自動化とかやりました。

興味というか次に調べたいのは *-module-source な
カーネルモジュールパッケージの作り方とかDKMSの使い方。
普通のパッケージと違う点が多々ありそう。
\end{prework}

\begin{prework}{ 野首 }

ホームのMHフォルダを外から見れるようuw-imapdを入れました。インストールし
 てdebconfに答えるだけでSSL readyなimap環境になりました。
\end{prework}

\begin{prework}{ MATOHARA }

2010年11月の勉強会資料を見ながらnilfs2 をバックアップディスクに設定して
 みました．
リサイズ機能も来たようなので試してみたいです．
- [PATCH 0/4] nilfs2 resize support -- Linux NILFS Development
\url{http://www.spinics.net/lists/linux-nilfs/msg00869.html}
\end{prework}

\begin{prework}{ hattorin }

バングラデシュでDebianインストールして、ネットワークの監視系ツールをいろ
 いろとセットアップしてきました。今はDebian on SqueezeでTokyoTyrantの
 KeyValueを使い、特定の問題を早く計算するためにネットワークを使って計算
 を早くするような仕組み作りをしています。
\end{prework}

\begin{prework}{ キタハラ }

実家でプリンタの設定したり、
スキャナーの設定に失敗してました。

\end{prework}

\begin{prework}{ koedoyoshida }

最近Debianで自分がやったこと
\begin{itemize}
\item ようやく、メイン環境をSqueezeにupdate
wide-dhcpがupdateできずにはまったが、下記を見て解決。
\url{http://www.flcl.org/~takasugi/tdiary-org/?date=20061023}
\item OSC仙台に参加、出展。
\end{itemize}


\end{prework}

\begin{prework}{ emasaka }

sidでちょっとはまったとこについて、GitHubでupstreamに簡単なパッチをpull
 requestしたら、そのバージョンが数日後にsidに降りてきました。
\end{prework}

\begin{prework}{ dictoss(杉本　典充) }

最近やっていることはkfreebsdを常用していること、Debianを開発環境として
 gtkアプリを試作していること。
今後はkfreebsdでIS03を使ってテザリング、klinuxよりkfreebsdの方がおすすめ
 といえる有利分野を見つけることをやりたい。
\end{prework}

\begin{prework}{ なかおけいすけ }

興味があること：DebianLive
最近ネカフェで一夜を明かしたのですが、ネカフェのPCにコンパイラが入ってな
 くて困ったので、USBやSDカードにインストールしたDebianを持ち歩いていると
 ハッピーになれるのではないかと。
\end{prework}

\begin{prework}{ 吉野(yy\_y\_ja\_jp) }

バグレポートとDDTSS/DDTPぐらいでしょうか．
\end{prework}

\begin{prework}{ henrich }

いくつかパッケージやメッセージをルーティンの更新しました。

\end{prework}

\begin{prework}{ Osamu MATSUMOTO }

\begin{itemize}
\item インフラ,サーバ管理の自動化\\
 自動インストール,構成管理,コンフィグ投入, 監視のまでを
 ラフに綺麗に繋ぎたい。Debian的な良い組み合わせあったら教えてください。
 (cobbler+ puppet/cfengine+ nagios + なんかwebcgi的な)
\end{itemize}

\end{prework}

\begin{prework}{ まえだこうへい }
\begin{itemize}
\item *diagシリーズ(http://blockdiag.com/)のdebパッケージ化中。
\item さくらのVPSを先月契約して、lxc \& Squeezeで開発\&検証環境に。
\item あらきさんメンテナンスのDebianのAMI使ってAWSでごにょごにょと。
\end{itemize}
\end{prework}

\begin{prework}{ yamamoto }

\begin{itemize}
\item 最近Debianで自分がやったこと\\
ポチポチと公式パッケージのリビルドをしました。
\item 興味のあること\\
移植。
\end{itemize}

\end{prework}

\begin{prework}{ 岩松 信洋 }

\begin{itemize}
\item SH4 buildd のメンテ。
\item libpng15 の experimental へのアップロードとtransition作業。
\item スポンサーのパッケージチェック、アップロードなど。
\end{itemize}

\end{prework}

\begin{prework}{ 野島　貴英 }

\begin{itemize}

\item pythonでgnome-notifyめがけて再生中のデータのメタデータ送る「できるだ
 け簡単にできる」totemのプラグイン書いてみた。
 \url{http://d.hatena.ne.jp/nozzy123nozzy/20110502/1304322969}
\item sidのalsa-lib,alsa-utils,alsa-driverの解析中。
（bluetoothヘッドフォン出力相手に、alsaのみで、演奏中の音をcaptureしたい
 し、音が出る複数のアプリの音をmixして出力したいっ）
\item sidのlinux-image-2.6.39-2-686-paeにて、multi-stateなusb装置にejectコ
 マンド発行すると高い確率でOopsする件のデバッグを誰かがパッチ出すまでや
 り中。（いったい誰だ、dptぶっ壊すのは...)

\end{itemize}

\end{prework}

\begin{prework}{ 上川純一 }
\begin{itemize}
\item xslt のツールチェインをいろいろといじってみたり。
\item Javascript のコードを書いてみたり。
\end{itemize}
\end{prework}

}

\emtext{XSLTPROC}

\begin{frame}{IRC会議システム}

 IRCで会議するときに議論が発散しないように議事録と議題をリアルタイムに表
 示するシステム

 \begin{itemize}
  \item 議事録はXML形式
  \item C++
  \item libdancer-xml
 \end{itemize}
\end{frame}

\begin{frame}{XSLT 1.0}
 XML形式で記述する文書処理言語
\end{frame}

\emtext{sphinx}
\emtext{今月のSuperH}
\emtext{月刊PPC64ポーティング}

\begin{frame}{今後のイベント}
 
\begin{itemize}
 \item 7月　Debian勉強会 \& Debconf11 in ボスニア
 \item 8月 未定
\end{itemize}
\end{frame}

\begin{frame}{今日の宴会場所}

参宮橋のどこか。

\end{frame}

\end{document}

;;; Local Variables: ***
;;; outline-regexp: "\\([ 	]*\\\\\\(documentstyle\\|documentclass\\|emtext\\|section\\|begin{frame}\\)\\*?[ 	]*[[{]\\|[]+\\)" ***
;;; End: ***
